\chapter{Intramedullary Nailing}

\section*{Introduction \& Clinical Indications}
Intramedullary nailing stabilizes a long-bone fracture with a metal rod placed within the medullary canal. It is commonly used for femoral, tibial, and humeral shaft fractures and selected metaphyseal or segmental injuries. Nails may be reamed or unreamed and locked with screws to control length, rotation, and alignment.

The method is chosen according to fracture pattern, soft-tissue injury, skeletal maturity, bone quality, infection, and associated injuries. External fixation, plating, casting, or staged fixation may be safer in selected open fractures, severe contamination, vascular injury, or unstable patients.

\section*{Relevant Surgical Anatomy}
Entry point, canal geometry, nearby nerves and vessels, fracture level, and joint surfaces guide nail selection. The femoral piriformis or trochanteric region, tibial proximal entry zone, and humeral entry structures have different risks. Rotation and length must be compared with the uninjured limb when possible.

\section*{Preoperative Workup \& Preparation}
Radiographs define the fracture and CT may be needed for complex extension. Neurovascular and compartment examinations are documented. Open fractures require antibiotics, tetanus care, and debridement. Consent includes infection, fat embolism, malalignment, nonunion, hardware failure, knee or shoulder pain, nerve injury, compartment syndrome, and reoperation.

\section*{Surgical Technique \& Procedure}
The patient is positioned for imaging and traction. The canal is entered through a safe approach, the fracture is reduced, and the nail is advanced across the fracture. Proximal and distal locking screws control rotation and length. Fluoroscopy confirms alignment, screw position, and joint clearance. Wounds are closed after hemostasis and compartment assessment.

\section*{Complications \& Risks}
Risks include infection, malrotation, shortening, nonunion, delayed union, implant breakage, iatrogenic fracture, neurovascular injury, fat embolism, compartment syndrome, knee or shoulder pain, and need for dynamization, exchange nailing, bone grafting, or revision.

\section*{Postoperative Care \& Recovery}
Weight-bearing and range of motion depend on fracture stability and associated injuries. Serial radiographs evaluate alignment and callus. Thromboprophylaxis, pain control, wound care, and physiotherapy are coordinated. Increasing pain, swelling, numbness, fever, drainage, or limb color change requires urgent review.

\section*{Outcomes \& Prognosis}
Intramedullary nails provide load-sharing stability and often permit early mobilization, but healing depends on fracture biology, alignment, soft tissue, smoking, diabetes, infection, and bone quality. There is no universal union or return-to-work rate across all bones and fracture patterns.

\section*{References}
\begin{enumerate}
  \item AO Foundation. Surgery Reference: Intramedullary Nailing.
  \item American Academy of Orthopaedic Surgeons. Fracture and Trauma Guidance.
  \item Orthopaedic Trauma Association. Principles of Fracture Care.
\end{enumerate}

\clearpage
